\section{Methodology}

\subsection{System Overview}
The proposed methodology consists of [number] main stages: [stage 1],
[stage 2], and [stage 3]. First, [describe the input acquisition or
preprocessing step]. Next, [describe the central algorithm, model, or
workflow]. Finally, [describe the output generation, decision process, or
postprocessing step].

\subsection{Data Preparation}
The study uses [name and description of the dataset or data source]. The
data are prepared by [cleaning, normalization, transformation, annotation,
or feature extraction procedure]. Any missing or invalid observations are
handled using [procedure], and the data are divided into [training,
validation, and test sets or other partitions] using [split strategy].

\subsection{Proposed Approach}
Let $x$ denote the input and let $y$ denote the corresponding target. The
proposed method learns a mapping
\begin{equation}
    \hat{y} = f(x;\theta),
\end{equation}
where $f$ is the proposed model and $\theta$ represents its trainable
parameters. The model is optimized using [loss function or optimization
objective]. The design choices are motivated by [briefly explain the
theoretical or practical rationale].

\subsection{Evaluation Criteria}
Performance is measured using [metric 1], [metric 2], and [metric 3]. These
metrics are selected because they measure [explain what each metric captures].
The proposed approach is compared with [baseline 1] and [baseline 2] under
the same evaluation conditions.
